\sectionguard
\section*{Appendix: Scope and Qualifications}

\editionnote

\subsection*{Edition, formulary identity and text-verification state}

The work is the Fifteenth Sunday after Pentecost of the \work{Missale Romanum ex
decreto Sacrosancti Concilii Tridentini restitutum, Summorum Pontificum cura
recognitum}, \latin{editio typica} (Vatican City: Typis Polyglottis Vaticanis,
1962). Printed heading \latin{DOMINICA / DECIMA QUINTA / post Pentecosten}, rank
\latin{II classis}, printed pp.~396--397, marginal nos.\ 1582--1591; \latin{Credo}
and the Preface of the Most Holy Trinity; no Tract, Sequence, second or third
oration, blessing, station or ritual text. Green follows RG 127~b of the 1960
\latin{Rubricae generales} and is not printed on these pages.

All ten elements, their headings, references, wording, accentuation, punctuation,
the printed rank, the rubrics, the marginal series and both formulary boundaries
were visually collated against the page images of the CMAA-hosted facsimile of
the Vatican typical edition on 2026-09-05, at the scan's native resolution, with
disputed glyphs enlarged. No transcription and no optical text layer was allowed
to supply a reading. The full record of that collation, its provenance and its
unresolved points is \texttt{propers/verified.md}; the unedited machine pull that
located the formulary is \texttt{propers/retrieved.txt}.

A second 1962 edition, the Benziger \latin{iuxta typicam}, was read only in an
uncorrected optical text layer. It is a different edition of the same revision and
not an independent image of the typical edition, so every difference between them
recorded in the commentary is an edition difference or an artefact of that layer.
Nothing was blended and no Benziger form was substituted for a Vatican typical
form. One word-form difference stands, \latin{Aduléscens} against
\latin{Adolescens} at Lk.\ 7:14.

\subsection*{Published text, English, and rights}

The basis for reproducing the Latin is 17 U.S.C.\ §103(b): the copyright in a
revised edition extends only to the material that edition contributed. Nine of
the ten elements are printed, in the same order and at the same references, in
three Missals tracked in this repository as public-domain text --- the Vatican
\latin{typica} of 1604, the Pustet Ratisbon 1862 and the Venice 1570 --- and the
Gospel in two of them, read in
uncorrected optical layers of which no page image was opened here. The source
records of the 1604 and the 1570 classify their layers degraded and forbid
resting a negative search on them; the Pustet's does neither. It is not a
page-image collation and
nothing is quoted from any of the three: where the commentary or the
appointed-text section reports one of their readings, a spelling or a point,
what it reports is what that layer shows, and the strings behind those reports
are the lightly normalised renderings recorded in \texttt{propers/verified.md}.
No negative result in this guide rests on them. The Pustet does not print this
Sunday's Gospel at all but cross-references it to the Lenten Thursday that
shares it, so the Gospel's presence in a public-domain witness rests on the 1604
and the 1570. The 1962 typical edition's own copyright status is unresolved and
is more likely than not restored in the United States to the end of 2057; the
reading of the statute above is this project's and is not a clearance anyone
granted, and it is not legal advice.

No English here is composed, translated, adapted or paraphrased. Scriptural
elements are quoted from the Douay--Rheims as revised by Bishop Challoner, in the
tracked Project Gutenberg e-text of that revision, with psalm loci resolved from
the tracked numbering concordance and not by arithmetic. The three orations are
quoted from \work{The Roman Missal translated into the English language for the
use of the laity}, first revised edition (Philadelphia: Eugene Cummiskey, 1861),
which this audit read on the page images of the registered scan at printed
pp.~427--429 and found word for word the tracked transcription. \textbf{That
English is a free devotional rendering and not a word-for-word translation}, its
translator is not established by any source opened here, and it is presented as an
identified historical English witness on that footing. Where the registered
English does not answer the Latin, the appointed-text section says so at the
element and leaves the gap open. At four sung elements the Douay has no
counterpart for a reading --- the Introit's added \latin{ad me}, the Alleluia's
\latin{super omnem terram}, the Offertory's three departed phrases, and the
Communion's \latin{dédero} and \latin{sǽculi} --- and beside them stand the
missal's own untranslated cues: the lesson's prefixed \latin{Fratres}, the
Gospel's incipit, and the Secret's conclusion, which the 1861 book leaves
untranslated. Of the four the Communion is a genuine negative
--- no English witness examined renders either reading distinctly.

\textbf{The repository's translation ledger and its own passage records
contradict each other about these three orations.} The ledger carries all three
as text-free rows with availability ``unavailable'' and a reason of
withheld rights, on the ground that no tracked publication artifact and passage
binding establishes text that may be served --- while that edition's tracked
public-domain artifact and three registered passages exist, two of them verified,
and this audit read all three orations on the registered page images. The English
printed here rests on the artifact and the passage records, which are cited; the
ledger is not cited as evidence either way, and it stands unreconciled with the
records it contradicts.

\subsection*{Evidence states and how claims are distinguished}

Verified source text is wording checked against an identified primary edition:
that is the Latin of all ten elements. A checked quotation is a published claim
read at an exact locus in the edition named at the point of use; that is every
patristic, medieval and liturgical citation here. Source-grounded synthesis is an
editorial conclusion supported by checked evidence: the claims of
\work{Source-Grounded Synthesis Across the Propers}. Editorial or AI proposal is
labelled as such and confined to \work{The Propers: Interpretive Possibilities}.
An unverified lead supports no published claim.

\textbf{Greek is given in transliteration throughout}, and the Greek witnesses are
the 1904 Nestle text for the New Testament and Swete's Septuagint for the psalms
and 3~Kingdoms; those, and no other Greek edition, stand behind every Greek word
printed. \textbf{No Greek patristic text was collated anywhere.} Cyril of
Alexandria on Luke and on John, Chrysostom on Galatians and on John, and Gregory
of Nyssa were all read in nineteenth-century English translation, and Cyril's Luke
commentary reaches English through an ancient Syriac version of a Greek original
surviving elsewhere only in catena fragments. Where one of those witnesses makes a
point about a Greek word, the point is reported as his argument and is not
established here.

\subsection*{Source scope, and what it does not reach}

Everything checked in Latin was read in Latin: Ambrose, Augustine, Jerome, Bede,
Hilary, Cassiodorus and Aquinas. The editions and their state vary and the
variation is material. Ambrose on Luke was read as an optical layer of the
critical edition, whose body text is sound and whose apparatus lines are mangled.
\textbf{Hilary on Ps.~91 was read as an optical layer of a scan of the critical
edition whose Latin recognition is poor}, whose running heads and apparatus are
badly mangled and whose page-number sequence is unreliable; CCSL~61/61A was not
consulted. Ambrose on Ps.~39 and Jerome on Galatians were read on page images with
no text layer and transcribed by eye. Bede and Cassiodorus were read on optical
layers of Migne whose recognition is poor and which corrupt individual words.
Augustine was read throughout in a web presentation of the Maurist text. Aquinas
was read in the standard scholarly electronic text of the Marietti editions.
\textbf{No critical edition later than 1919 was consulted anywhere}: those are in
copyright and are not held here.

Named commentators the repository's own index ranks for these passages and which
were not opened: on the four psalms, Theodoret, Jerome, Arnobius the Younger,
Bellarmine, Peter Lombard, Nicholas of Lyra, Euthymius Zigabenus, Bruno the
Carthusian, Albert the Great and Hugh of Saint-Cher; on the Epistle,
Ambrosiaster, Theodoret, Marius Victorinus, Theophylact, Haymo of Halberstadt, the
\work{Glossa ordinaria}, Denis the Carthusian, Cornelius a~Lapide and William
Estius; on the Gospel, Theophylact, Titus of Bostra, Euthymius Zigabenus, Albert
the Great, Bonaventure, Hugh of Saint-Cher and Maldonado; on the Communion's
verse, Hilary's \work{De Trinitate} VIII, Irenaeus, Theophylact, Rupert of Deutz
and Maldonado. Bonaventure's exposition of the Communion's verse was located in a
tracked public-domain volume and not read.

On the liturgical side: Rheinau, S.~Gallen 348, Pamelius, Menard and Gerbert's
printed p.~175 were not opened, so the five-witness Sunday-numbering split is
Wilson's report of those books. The Wurzburg comes and Klauser's \work{Capitulare
evangeliorum} were not sought, so the lectionary result is what Migne prints at
PL~30 under Jerome's name, whose date, provenance and authority nobody here
established. No Ordo Romanus, station list, Breviary of any date, Ambrosian or
Mozarabic missal was opened for this Sunday, and whether it has a station church
is not claimed either way. \textbf{Hesbert's \work{Antiphonale Missarum
Sextuplex} cannot be reached under the rights policy in force}: the edition is
registered here at catalogue level with no artifact, on the ground that no
affirmative United States redistribution basis for exact edition bytes was
established, so every antiphonary datum in this guide is the gregorien.info
database's report of his pages and is attributed to that database.

\subsection*{Search limits and material negative results}

Every negative below is the answer one corpus gave to one query, and a different
corpus, a different phrasing or a non-English literature could overturn any of
them.

\begin{itemize}
\item No Bible witness and no patristic or Greek lemma collated here reads
\latin{super omnem terram} at Ps.\ 94:3 or \latin{miserere mihi} at Ps.\ 85:3.
Both readings stand in the optical layers of the three older Missals read for
the preexisting-material check, which are neither Bible text nor patristic
lemma. Nothing in Swete answers the Introit's added \latin{ad me}, which Cassiodorus's
Latin lemma nonetheless carries.
\item The Communion's \latin{dédero} and \latin{sǽculi} answer nothing in the
Greek, and no English witness examined --- the tracked Douay, the Keating of 1806,
the Cummiskey of 1861, the Lasance of 1945, Guéranger's 1909 volume --- renders
either distinctly.
\item \latin{Naim} occurs once in the whole tracked Clementine, and \latin{propheta
magnus} only at Ecclus.\ 48:25 and Lk.\ 7:16.
\item The eight-resurrections scheme the \work{Catena aurea} attributes to
Ambrose is not in his Luke commentary, on a word search of the complete critical
text; it may stand in another of his works, or be Ps.-Ambrose, or be a Glossa
attribution.
\item Gregory the Great has nothing on this pericope, on the evidence of the
\work{Catena}'s lemma frequencies over Luke~7 and of the standard account of his
forty homilies. PL~76 and CCSL~141 were not opened.
\item No Father or Doctor comments on the three orations anywhere in this
repository's tracked sources or in the patristic and medieval works opened for
this formulary; the patristic corpora at large were not searched, and a Father
found quoting one of the three incipits corrects the result. None of the three
occurs in the Veronense in two independent scans, in Férotin's \work{Liber
mozarabicus sacramentorum}, in Bannister's \work{Missale Gothicum} or in Lowe's
\work{Bobbio Missal}. \textbf{Four bounds travel with that second negative}: it is
a literal search over uncorrected optical text and cannot see a paraphrase or a
recentonised descendant; three of the four books are diplomatic editions printing
manuscript contractions; the \latin{e caudata} failure demonstrated below would
hide any of the three in any of the four books; and the corpus itself excludes
every eighth-century Gelasian, every Ambrosian book, every Gregorian supplement
beyond Wilson, every medieval monastic or diocesan sacramentary and every later
Roman missal.
\item No registered passage record exists anywhere in this repository's source
library at any locus of Luke~7, at any of the four appointed psalms, or at any
verse of the appointed Epistle; and no Graduale, Antiphonale, Liber usualis,
Cantatorium or Latin psalter other than the Clementine is registered at all. No
Psalterium Romanum, no Psalterium \latin{iuxta Hebraeos} and no Old Latin psalter
was collated for this formulary.
\item \textbf{No qualifying cultural afterlife was located for the Introit, the
Gradual, the Alleluia or the Communion}, across four corpora searched on
2026-09-05: the Library of Congress Chronicling America collection, the
CourtListener opinion index, the UK Parliament Hansard contributions index and the
Caselaw Access Project static corpus. The counts are recorded in
\texttt{research/scope.md}. The Alleluia is closed to that sweep as a matter of
text and not of effort, because no English corpus can be searched for a reading
the English Bibles do not carry. \textbf{On this selection the gallery has no
Gospel entry}: the one located non-preaching use of the appointed Gospel is
contestable under the profile's own rule, and the absence is left standing.
\item Those four corpora are overwhelmingly American and English-language. No
British, Irish, Canadian or Australian newspaper corpus, no non-English
literature, no film or broadcast archive beyond two catalogue records, no
advertising archive and no manuscript source was searched.
\item The precedent search behind the five interpretive proposals ran inside
this repository alone: 669 proper-guide files across both providers and both
calendars, 628 article files, both providers' theology and history trees, the
1962 Ordinary and reference works beside this series, and the tracked 1962
calendar registry and commentary corpus index. The open web, external sermon
and homiletic databases, printed chant repertories and any universal
literature search lie outside it, so \emph{not located in the checked corpus}
says only that this bounded corpus does not hold the conjunction.
\item No gallery entry rests on a search snippet: every legal candidate retained
was read in the Caselaw Access Project's continuous text of the printed reporter
and the phrase confirmed there. \textbf{No page-image collation stands behind any
gallery entry}, so none of the five is described as verified and the printed
pointing of none of them is presented as settled.
\end{itemize}

\textbf{One method warning governs every literal search reported here.} A search
over registered optical text layers produces false negatives, and the
demonstration is on this formulary's own prayers: a first pass over the Hadrianum
layer returned nothing for \latin{miseratio continuata}, \latin{Mentes nostras et
corpora} and \latin{doni caelestis operatio}, and a second pass found all three,
the layer rendering a footnote marker inside one word and a caret for the
\latin{ae}-ligature in another. Every negative here was run with loosened patterns
and cross-checked against a second scan or a second pattern where one existed, and
each remains correctable.

\subsection*{Chronology}

Every biblical date printed in this work comes from the Scripture chronology
corpus and from nowhere else, carried in \texttt{research/chronology.toml} and
projected mechanically into the Date column of the page-2 sheet. No date here was
inferred, harmonised, or taken from a commentary or a chronological table. For all
four psalm elements the composition answer is a single whole-Psalter boundary on
a modern critical profile, which states in its own words that no individual
psalm can be dated securely. For the three psalms whose titles name David the
corpus also answers a traditional-attribution relation, marked disputed, which
carries David's reign in the usual chronology as a named reference point and not
as a date of the psalm; that relation leads the Date cell, the explanatory row
names its source, and the Gradual, whose title names no author, carries none.
The Epistle
carries five labels spanning nine years within one traditional profile, one
preferred and four alternate, and all five are printed. The Gospel and the
Communion each carry two disputed composition labels and one preferred
narrated-event label, and the two kinds answer different questions --- when the
book was written and when the thing it tells of happened.

\subsection*{Review state}

The research behind this guide was carried out in one production and its audit
record is \texttt{research/scope.md}, which holds the passage-by-passage reception
matrix, the corpora and editions searched, the rejected and unresolved leads, the
competing historical judgments, the gallery audit with its candidates not
selected, and the audit of the interpretive proposals with their precedent
classifications and search boundary. The bindings this publication makes to the
repository's reusable source library are recorded in
\texttt{research/source-bindings.toml}. This is a hand-missal companion for study,
not an official liturgical text, a critical edition or a substitute for the cited
sources, and it has passed no specialist, clerical or ecclesiastical review.
